\chapter{Problems and Lessons Learned}
\label{chap:lessons}

Every project like this accumulates real problems along the way: bad data, bugs, and results
that do not go the way you expected. This chapter collects the ones from this project
honestly, rather than smoothing them out of the narrative.

\begin{enumerate}[leftmargin=1.4em]
  \item \textbf{No absolute clock in the raw files.} The recordings and their annotations each
    started their own internal clock at zero, with no shared reference point anywhere in the raw
    data. Recovering the real offset (\Cref{chap:rawdata}) required rejecting two tempting
    shortcuts that would have quietly fabricated seizure labels.
  \item \textbf{Conflicting annotation rows.} Two recordings had genuinely contradictory seizure
    times in the workbook. Rather than pick a version by hand, each candidate was scored against
    the real EEG signal, and the better-separated candidate was kept --- a repeatable, checkable
    decision instead of a guess.
  \item \textbf{Class imbalance by construction.} Only 3.96\% of windows contain a seizure.
    Balanced training and the first augmentation attempt (shift+mirror) both turned out to
    \emph{hurt} overall event F1 despite improving recall --- a genuine negative result, reported
    in full rather than hidden. The lesson was not to stop there: testing five more augmentation
    methods individually (\Cref{chap:augmentation}) found two, channel dropout and noise, that
    genuinely help, and channel dropout became this project's best result. A negative result for
    the first method tried was not a verdict on the whole approach.
  \item \textbf{A real crash in the CNN training script.} \texttt{event\_counts()} used a pandas
    index label as if it were an array position, which worked by coincidence until a particular
    random validation split exposed it. The bug, how it was found, and the fix are described in
    full in \Cref{chap:deeplearning}.
  \item \textbf{Deep learning did not beat the classical baseline.} This was treated as a
    hypothesis to test, not a result to engineer toward. It did not beat the SVM, and that
    outcome is reported as the finding it is, in \Cref{chap:deeplearning}, rather than reframed
    as a near-win.
  \item \textbf{A long evaluation run had to be killed.} The nested-threshold re-evaluation of
    the classical models (\Cref{chap:evaluation}) is memory-heavy and was killed once for lack of
    available RAM during a full attempt; only 3 models have been scored this way so far. This
    remains an open limitation, not something to paper over, and is listed again in
    \Cref{chap:conclusion}.
  \item \textbf{A naming collision silently overwrote good results.} When bringing new CNN runs
    back from the GPU machine, a quick, one-fold smoke test (meant only to check the training
    script still worked after a code change) happened to reuse the exact output names of two
    already-trained, full, five-fold models. Copying the smoke test's files into place overwrote
    the real results with a tiny, partial run until the mismatch was caught --- the overwritten
    run's window count and found-seizure count were obviously too small for a full run, which is
    what gave it away. The fix was mechanical (restore the real files, re-copy only the genuinely
    new runs), but the lesson is not: any script whose output name depends on its flags needs a
    naming scheme where a debug run and a real run can never collide, and every number brought
    back from an external run is worth a sanity check against what a full run should look like
    before it is trusted.
\end{enumerate}

\begin{summarybox}
Nothing in this project went perfectly on the first try. The raw data needed real detective
work before it could be trusted; two of the three attempts to fix class imbalance made results
worse, not better; a real bug crashed the CNN training script on certain random splits; and a
long-running evaluation script had to be killed for memory. All of it is reported here rather
than quietly cleaned up, because a negative result, honestly reported, is still useful
information.
\end{summarybox}
